\begin{description}
\item[a)]
  \begin{equation*}
  \frac{1}{2}m_1 v_1^2 + \frac{1}{2}m_2 v_2^2 - \frac{Gm_1 m_2}{r_1 + r_2} = E.
  \tag*{[5]}
  \end{equation*}
  Note that the total energy is a negative quantity.

\item[b)] From $m_1 r_1 = m_2 r_2$ and $r_1 + r_2 = r$, we have
  \[
  r_1 = \frac{m_2}{m_1 + m_2}\,r
  \qquad\text{and}\qquad
  r_2 = \frac{m_1}{m_1 + m_2}\,r.
  \]
  Substituting them into the solution obtained in a), we get
  \begin{align*}
  \frac{1}{2}\mu v^2 - \frac{G(m_1 + m_2)\mu}{r} &= E \\
  v^2 - \frac{2GM}{r} &= \frac{2E}{\mu}
  \end{align*}
  Therefore,
  \begin{equation*}
  v^2 = 2\left(\frac{GM}{r} + \frac{E}{\mu}\right)
  \tag*{[10]}
  \end{equation*}

\item[c)] From b), we have
  \begin{align*}
  v^2 &= \frac{2GM}{r} + \frac{2E}{\mu}
  = (2GM)\left(\frac{1}{r} + \frac{E}{GM\mu}\right) \\
  &= (2GM)\left(\frac{1}{r} - \frac{1}{r_0}\right),
  \quad\text{where } r_0 = -\frac{GM\mu}{E}.
  \tag*{[5]}
  \end{align*}

\item[d)] Substituting $r(\theta)$ in the equation of part c):
  \[
  v^2 = 2GM\left(\frac{1}{r} - \frac{1}{r_0}\right)
  = 2GM\left(\frac{1}{\tfrac{1}{2}r_0(1-\cos\theta)} - \frac{1}{r_0}\right)
  = \frac{2GM}{r_0}\left(\frac{1+\cos\theta}{1-\cos\theta}\right)
  \]
  \[
  v^2 = \frac{2GM}{r_0}\,\frac{(1-\cos^2\theta)}{(1-\cos\theta)^2}
  = \frac{2GM}{r_0}\,\frac{\sin^2\theta}{(1-\cos\theta)^2}
  \;\Longrightarrow\;
  v(\theta) = \left(\frac{2GM}{r_0}\right)^{1/2}\frac{\sin\theta}{1-\cos\theta}
  \]
  Combining $v(\theta)$, $t(\theta)$, and $r(\theta)$:
  \[
  \frac{v(\theta)\,t(\theta)}{r(\theta)}
  = \frac{\left(\dfrac{2GM}{r_0}\right)^{1/2}
      \dfrac{\sin\theta}{1-\cos\theta}\cdot
      \left(\dfrac{r_0^3}{8GM}\right)^{1/2}(\theta - \sin\theta)}
      {\dfrac{r_0}{2}\,(1-\cos\theta)}
  = \frac{2}{r_0}\left(\frac{2GM r_0^3}{r_0\, 8GM}\right)^{1/2}
    \frac{(\sin\theta)(\theta - \sin\theta)}{(1-\cos\theta)^2}
  \]
  \[
  \frac{vt}{r} = \frac{(\sin\theta)(\theta - \sin\theta)}{(1-\cos\theta)^2}
  \]

  \emph{Alternative solution:}
  From $r(\theta) = \dfrac{r_0}{2}(1-\cos\theta)$, we have
  \[
  v = \left(\frac{r_0}{2}\sin\theta\right)\dot\theta
  \]
  From $t(\theta) = \left(\dfrac{r_0^3}{8GM}\right)^{1/2}(\theta-\sin\theta)$,
  we have
  \[
  1 = \left(\frac{r_0^3}{8GM}\right)^{1/2}(1-\cos\theta)\,\dot\theta
  \]
  \begin{equation*}
  \therefore\;
  \frac{vt}{r}
  = \frac{\left(\dfrac{r_0}{2}\sin\theta\right)
      \left(\dfrac{r_0^3}{8GM}\right)^{1/2}(\theta-\sin\theta)}
      {\left(\dfrac{r_0^3}{8GM}\right)^{1/2}(1-\cos\theta)\,
       \dfrac{r_0}{2}(1-\cos\theta)}
  = \frac{(\sin\theta)(\theta-\sin\theta)}{(1-\cos\theta)^2}.
  \tag*{[10]}
  \end{equation*}

\item[e)]
  \begin{align*}
  v &= -118~\mathrm{km\,s^{-1}} = -118 \times 10^{3}~\mathrm{m\,s^{-1}} \\
  r &= 710~\mathrm{kpc} = 2.19 \times 10^{22}~\mathrm{m} \\
  t &= 13700 \times 10^{6}~\mathrm{years} = 4.3233 \times 10^{17}~\mathrm{s}
  \end{align*}
  \begin{equation*}
  \therefore\;
  \frac{vt}{r}
  = -\frac{118 \times 10^{3} \times 4.3233 \times 10^{17}}
          {2.1908 \times 10^{22}}
  = -2.33
  = \frac{(\sin\theta)(\theta-\sin\theta)}{(1-\cos\theta)^2}
  \tag*{[5]}
  \end{equation*}
  The negative value of the left-hand side of this equation implies that
  $\theta$ is greater than $\pi$.

  \begin{center}
  \begin{tabular}{cc}
  $\theta$ & $(\sin\theta)(\theta-\sin\theta)/(1-\cos\theta)^2$ \\
  \hline
  $200^\circ = 3.49$ radians & $-0.348$ \\
  $210^\circ = 3.67$ radians & $-0.598$ \\
  $240^\circ = 4.19$ radians & $-1.95$ \\
  $244^\circ = 4.26$ radians & $-2.24$ \\
  $245^\circ = 4.28$ radians & $-2.32$ \\
  $246^\circ = 4.29$ radians & $-2.40$ \\
  $250^\circ = 4.36$ radians & $-2.77$ \\
  \end{tabular}
  \end{center}
  This gives the value of $\theta = 245^\circ = 4.28$ radians. \hfill[5]

\item[f)] From $r(\theta) = \dfrac{r_0}{2}(1-\cos\theta)$, we have
  $r_{\mathrm{max}} = r_0$, and hence
  \begin{equation*}
  r_{\mathrm{max}} = \frac{2r(\theta)}{1-\cos\theta}
  = \frac{2 \times 710~\mathrm{kpc}}{1 - \cos 245^\circ}
  = 998~\mathrm{kpc}
  \tag*{[5]}
  \end{equation*}
  (This is the maximum separation between the Milky-Way Galaxy and Andromeda
  Galaxy; they started to move towards each other afterwards.)

  The value of $M$ can be calculated from the relation
  $t(\theta) = \left(\dfrac{r_0^3}{8GM}\right)^{1/2}(\theta - \sin\theta)$:
  \[
  M = \left(\frac{r_{\mathrm{max}}^3}{8G}\right)
      \left\{\frac{\theta - \sin\theta}{t(\theta)}\right\}^{2}.
  \]
  From the maximum separation of
  $r_{\mathrm{max}} = 998~\mathrm{kpc} = 3.0795 \times 10^{22}$~m,
  $\theta = 245^\circ = 4.276$~radians,
  $t = t(\theta) = 13700 \times 10^{6}~\mathrm{years}
     = 4.3233 \times 10^{17}$~s, and
  $G = 6.67 \times 10^{-11}~\mathrm{N\,m^2\,kg^{-2}}$, we obtain
  \begin{equation*}
  \therefore\; M = 7.86 \times 10^{42}~\mathrm{kg}
  = \frac{7.86 \times 10^{42}}{1.99 \times 10^{30}}\,M_\odot
  = 3.95 \times 10^{12}\,M_\odot
  \tag*{[5]}
  \end{equation*}
  The estimated number of stars, most of which are of less than a solar mass,
  for Andromeda and Our Galaxy are $4 \times 10^{11}$ and $2 \times 10^{11}$
  respectively. This implies that most of mass of the system is DARK.
\end{description}

\medskip
\noindent\textbf{Note that}
\[
v = \frac{d}{dt}r = +(2GM)^{1/2}\left(\frac{r_0 - r}{r_0 r}\right)^{1/2}
\quad\text{for } r_0 > r
\]
\[
\int (2GM)^{1/2}\,dt = +\int \left(\frac{r_0 r}{r_0 - r}\right)^{1/2} dr
\]
Putting $\dfrac{r}{r_0} \equiv \sin^2\phi$,
\[
t = \left(\frac{r_0^3}{8GM}\right)^{1/2}\{2\phi - \sin 2\phi\} + D
\]
We choose the initial condition such that $m_1$ and $m_2$ are close together,
that is that $r = 0$ at $t = 0$. This implies that $\phi = 0$ at $t = 0$,
hence $D$ must be zero.
\[
t = \left(\frac{r_0^3}{8GM}\right)^{1/2}(\theta - \sin\theta),
\quad \theta \equiv 2\phi.
\]
And from $r = r_0 \sin^2\phi = \dfrac{1}{2}r_0(1 - \cos 2\phi)$, we have:
\[
r = \frac{1}{2}r_0(1 - \cos\theta).
\]
These form the parametric solution:
$r_{\mathrm{max}} = r(\theta = \pi) = r_0$.
